(1) 质点坐标 \((s\sin\theta\cos\Omega t,s\sin\theta\sin\Omega t,s\cos\theta)\)，故 \(v^2=\dot s^2+\Omega^2s^2\sin^2\theta\)。拉氏量 \(L=\frac m2(\dot s^2+\Omega^2s^2\sin^2\theta)-mgs\cos\theta\)；Euler–Lagrange 方程为 \(\ddot s=\Omega^2s\sin^2\theta-g\cos\theta\)。
(2) 记 \(\kappa=\Omega\sin\theta\)，在 \(\Omega>0\) 时平衡位置 \(s_*=g\cos\theta/(\Omega^2\sin^2\theta)>0\)。约化有效势 \(U_{\rm eff}=mgs\cos\theta-\frac12m\kappa^2s^2\)，故 \(U_{\rm eff}''=-m\kappa^2<0\)，是局部最大值，平衡在线性层面严格不稳定而非稳定离心悬浮。
(3) 令 \(y=s-s_*\)，则 \(\ddot y-\kappa^2y=0\)，初值解为 \(s(t)=s_*+(s_0-s_*)\cosh(\kappa t)+(v_0/\kappa)\sinh(\kappa t)\)。一般扰动有 \(e^{\kappa t}\) 增长，增长率 \(\kappa=\Omega\sin\theta\)；若珠子到达槽端点 \(s=0\) 或有限槽长，则应停止该自由槽模型。
(4) 槽带动珠子绕轴所需的力矩为 \(\tau_z=d(m\Omega s^2\sin^2\theta)/dt=2m\Omega s\dot s\sin^2\theta\)，功率 \(P_{\rm drive}=\tau_z\Omega=2m\Omega^2s\dot s\sin^2\theta\)。实验室机械能 \(E_{\rm lab}=\frac m2(\dot s^2+\kappa^2s^2)+mgs\cos\theta\) 满足 \(\dot E_{\rm lab}=P_{\rm drive}\)；相反 \(J=\frac m2\dot s^2+U_{\rm eff}(s)\) 的导数为 \(m\dot s[\ddot s-\kappa^2s+g\cos\theta]=0\)。
